%% kapa: exam=2 round=25 year=2014 period=3 subject=compensation_law subject_ko=감정평가 및 보상법규 seq=3 number=3 points=20
%% source: Q-Net 감정평가사 기출문제 (cst00309) artlSeq=5109209 file=extracted/1008/5109209/2014년 제25회 감정평가사 2차시험 기출문제/3교시(감정평가 및 보상법규).hwp
%% conversion: kordoc 4.24.51 -> tools/gen2.py (auto); review=proofread-sonnet; flags=none
\begin{longquestion}{【문제 3】}
법원으로부터 근저당권에 근거한 경매를 위한 감정평가를 의뢰받은 감정평가사 乙이 감정평가 대상토지의 착오로 실제 대상토지의 가치보다 지나치게 낮게 감정평가액을 산정하였다. 토지소유자인 甲이 이에 대해 이의를 제기하였음에도 경매담당 법관 K는 乙의 감정평가액을 최저입찰가격으로 정하여 경매절차를 진행하였으며, 대상토지는 원래의 가치보다 결국 낮게 丙에게 낙찰되어 甲은 손해를 입게 되었다. 甲이 법관의 과실을 이유로 국가배상을 청구한 경우 이 청구의 인용가능성을 검토하시오. (20점) \par
\end{longquestion}
